\documentclass[11pt]{article}
\usepackage{amsmath}
\usepackage{amssymb}
\usepackage{algorithm}
\usepackage{algorithmic}
\usepackage{geometry}
\geometry{margin=1in}

\title{Step-Size Parameter Computation for PDHG Algorithm\\
with TV and L1 Regularization\\
\large Response to Dr. Sidky's Questions}
\author{Technical Analysis}
\date{\today}

\begin{document}

\maketitle

\section{Executive Summary}

This document provides detailed answers to Dr. Sidky's questions regarding the DBT reconstruction implementation:

\begin{enumerate}
    \item \textbf{ROI Visualization}: The dark background in reconstructed images is due to the ROI being extracted from the breast tissue region of the VICTRE phantom, with surrounding air appearing black (value 0).
    
    \item \textbf{Projection Geometry}: The implementation uses a \textbf{stack of 2D fan-beam projections}, not true 3D cone-beam projection. Each z-slice is processed independently with its own 2D fan-beam forward/backward projection.
    
    \item \textbf{Step-Size Parameters}: Detailed computation provided below for the PDHG algorithm with combined TV and L1 regularization.
\end{enumerate}

\section{Step-Size Parameter Computation Details}

\subsection{Optimization Problem Formulation}

The L1-DTV reconstruction solves the following optimization problem:
\begin{equation}
\min_{x \geq 0} \frac{1}{2}\|Ax - b\|_2^2 + \alpha\|\nabla_x x\|_1 + \beta\|\nabla_y x\|_1 + \gamma\|x\|_1
\end{equation}

where:
\begin{itemize}
    \item $A$: Fan-beam forward projection operator
    \item $b$: Measured sinogram data
    \item $x$: Image to reconstruct (constrained to be non-negative)
    \item $\nabla_x, \nabla_y$: Directional gradient operators
    \item $\alpha = 2 - 1.75 = 0.25$: Weight for x-direction gradients
    \item $\beta = 1.75$: Weight for y-direction gradients
    \item $\gamma = 5.0$: L1 sparsity regularization weight
\end{itemize}

\subsection{Primal-Dual Formulation}

The problem is reformulated as a saddle-point problem:
\begin{equation}
\min_{x \geq 0} \max_{y} \langle Kx, y \rangle - g^*(y) - f(x)
\end{equation}

where the combined linear operator is:
\begin{equation}
K = \begin{bmatrix}
A \\
\nabla_x \\
\nabla_y \\
I
\end{bmatrix}
\end{equation}

\subsection{PDHG Algorithm with He-Yuan Predictor-Corrector}

The algorithm alternates between primal and dual updates:

\begin{algorithm}
\caption{PDHG with He-Yuan Predictor-Corrector}
\begin{algorithmic}
\STATE \textbf{Initialization:} $x^0 = 0$, $y^0 = 0$, compute $\tau$, $\sigma$
\FOR{$k = 0$ to $\text{maxiter}$}
    \STATE \textbf{Primal Update:}
    \STATE $x^{k+1} = \text{proj}_{x \geq 0}\left(x^k - \tau K^T y^k\right)$
    \STATE $\bar{x}^{k+1} = x^{k+1} + (x^{k+1} - x^k)$ \COMMENT{Over-relaxation}
    
    \STATE \textbf{Dual Update:}
    \STATE $\tilde{y}^{k+1} = y^k + \sigma K\bar{x}^{k+1}$
    \STATE Apply proximal operators for each dual variable
    
    \STATE \textbf{Predictor-Corrector:}
    \STATE $y^{k+1} = y^k - \rho(y^k - \tilde{y}^{k+1})$ where $\rho = 1.75$
\ENDFOR
\end{algorithmic}
\end{algorithm}

\subsection{Operator Norm Estimation}

The operator norm $L = \|K\|$ is estimated using power iteration:

\begin{algorithm}
\caption{Operator Norm Estimation via Power Iteration}
\begin{algorithmic}
\STATE Initialize random $x \sim \mathcal{N}(0, 1)$
\FOR{200 iterations}
    \STATE \textbf{Forward (apply $K$):}
    \STATE $w = A(x) \cdot \nu_{\text{sino}}$ \COMMENT{Scaled projection}
    \STATE $g_x = \nabla_x(x) \cdot \nu_x$ \COMMENT{Scaled x-gradient}
    \STATE $g_y = \nabla_y(x) \cdot \nu_y$ \COMMENT{Scaled y-gradient}
    \STATE $g_{L1} = \gamma \cdot x$ \COMMENT{L1 term}
    \STATE $\text{mag}_1 = \sqrt{\|w\|_2^2 + \|g_x\|_2^2 + \|g_y\|_2^2 + \|g_{L1}\|_2^2}$
    
    \STATE \textbf{Backward (apply $K^T$):}
    \STATE $x_1 = A^T(w) \cdot \nu_{\text{sino}}$
    \STATE $x_2 = \text{div}_x(g_x) \cdot \nu_x$
    \STATE $x_3 = \text{div}_y(g_y) \cdot \nu_y$
    \STATE $x = x_1 + x_2 + x_3 + \gamma \cdot g_{L1}$
    \STATE $\text{mag}_2 = \|x\|_2$
    
    \STATE Normalize: $x \leftarrow x/\text{mag}_2$
\ENDFOR
\STATE \textbf{Return:} $L = (\text{mag}_1 + \text{mag}_2)/2$
\end{algorithmic}
\end{algorithm}

\subsection{Step-Size Formulas}

Given the operator norm $L$ and a balancing parameter $\text{stepbalance} = 100$:

\begin{align}
\tau &= \frac{1}{L \cdot \text{stepbalance}} \\
\sigma &= \frac{\text{stepbalance}}{L}
\end{align}

For the two-channel variant:
\begin{align}
\sigma_{\text{high}} &= \sigma \\
\sigma_{\text{low}} &= 4.0 \cdot \sigma \quad \text{(allowing more aggressive low-frequency denoising)}
\end{align}

\subsection{Convergence Guarantee}

The algorithm converges when:
\begin{equation}
\tau \cdot \sigma \cdot L^2 < 1
\end{equation}

With our choice:
\begin{equation}
\tau \cdot \sigma = \frac{1}{L \cdot \text{stepbalance}} \cdot \frac{\text{stepbalance}}{L} = \frac{1}{L^2}
\end{equation}

Therefore, $\tau \cdot \sigma \cdot L^2 = 1$, which is at the boundary of the convergence region. The predictor-corrector step with $\rho = 1.75$ provides additional stability.

\section{Implementation Details}

\subsection{Scaling Factors}

The implementation uses pre-computed scaling factors:
\begin{itemize}
    \item $\nu_{\text{sino}} = 1/\|A\|$: Projection operator scaling
    \item $\nu_x = 0.5/\|\nabla_x\|$: X-gradient scaling (nuxfact = 0.5)
    \item $\nu_y = 0.5/\|\nabla_y\|$: Y-gradient scaling (nuyfact = 0.5)
\end{itemize}

\subsection{Parameter Values}

\begin{table}[h]
\centering
\begin{tabular}{|l|c|l|}
\hline
\textbf{Parameter} & \textbf{Value} & \textbf{Description} \\
\hline
$\alpha$ & 1.75 & DTV parameter (controls directional weighting) \\
$\beta$ & 5.0 & L1 weight \\
$\rho$ & 1.75 & Predictor-corrector parameter \\
$\epsilon$ & 0.001 & Data discrepancy tolerance (RMSE) \\
stepbalance & 100.0 & Step-size balancing factor \\
cutoff\_lo & 8.0 & Low-frequency cutoff (two-channel) \\
$\sigma_{\text{lo}}$ scale & 4.0 & Low-frequency step-size multiplier \\
\hline
\end{tabular}
\end{table}

\section{Key Observations}

\begin{enumerate}
    \item \textbf{Fan-beam approximation}: The slice-by-slice fan-beam approach is computationally efficient and adequate for DBT's limited-angle geometry where cone-beam effects are minimal.
    
    \item \textbf{Step-size balance}: The large stepbalance value (100) prioritizes dual variable updates, which is appropriate for the regularization-heavy nature of the problem.
    
    \item \textbf{Two-channel advantage}: The frequency-split approach with different step sizes allows more aggressive denoising of low frequencies while preserving high-frequency details.
    
    \item \textbf{Predictor-corrector stabilization}: The He-Yuan scheme with $\rho = 1.75$ provides additional convergence stability beyond the basic PDHG algorithm.
\end{enumerate}

\section{Recommendations for Response to Dr. Sidky}

\begin{enumerate}
    \item Clarify that the dark background is expected behavior - the ROI focuses on clinically relevant breast tissue.
    
    \item Acknowledge the fan-beam approximation and explain it's a common simplification in DBT reconstruction that provides good results with computational efficiency.
    
    \item Emphasize that the step-size computation follows standard PDHG theory with careful operator norm estimation to ensure convergence.
    
    \item Consider mentioning that true cone-beam reconstruction could be implemented but would increase computational cost with minimal benefit for the thin-slice DBT geometry.
\end{enumerate}

\end{document}